\documentclass[11pt,a4paper]{article}
\usepackage[T1]{fontenc}
\usepackage{lmodern}
\usepackage[margin=25mm]{geometry}
\usepackage{amsmath,amssymb,amsthm}
\usepackage{booktabs,longtable}
\usepackage[hidelinks]{hyperref}
\usepackage{url}
\newtheorem{theorem}{Theorem}
\newtheorem{lemma}{Lemma}
\newcommand{\R}{\mathcal R}
\newcommand{\Q}{\mathcal Q}
\newcommand{\F}{\mathcal F}
\newcommand{\ind}{\mathbf 1}
\setlength{\emergencystretch}{3em}
\title{An explicit optimality certificate for\\
\texttt{r240\_p120\_m120\_s5.txt}}
\author{}
\date{4 October 2026}

\begin{document}
\maketitle

\begin{abstract}
For the specified instance and the passenger--parcel routing rules stated
below, the maximum total benefit is exactly $13\,279$. We prove the upper
bound by assigning nonnegative integer weights to taxis, requests, and four
three-resource conflict groups. These weights cover the benefit of every
feasible route. A feasible fleet can use each weighted resource or conflict
group at most once, so its benefit cannot exceed the sum of the weights,
$10\,563+2\,547+169=13\,279$. An explicitly supplied fleet attains this value.
The numerical certificate can be checked with integer arithmetic without
an optimization solver. Its completeness depends on the exhaustive route
enumeration justified in this document, rather than on ALNS convergence.
\end{abstract}

\section{Precisely which problem is being proved}

The input file is \path{data/probation/r240_p120_m120_s5.txt}. It has $20$
taxis, $120$ passenger requests, and $120$ parcel requests. Its SHA-256 is
recorded in Appendix~\ref{sec:reproduce}; the claim concerns those exact bytes.
The routing semantics follow \path{docs/problem/problem.tex} and the current full
solution evaluator.

A taxi starts at its own origin at time zero. Requests are optional. A served
request has its pickup and delivery performed once, in that order, by the
same taxi. A passenger's delivery must immediately follow its pickup;
parcels may remain onboard during this passenger trip. Parcel load is never
greater than the taxi's capacity. The stated service duration applies at
both endpoints. Passenger pickups and both parcel endpoints have time
windows: arrival must be no later than the latest time, and an early arrival
waits until the earliest time before service. Passenger deliveries have no
time window. A route ends after its last service point, without a return leg.
Travel times and costs are the input's ordered-pair values; no symmetry or
triangle inequality is assumed.

Let $P_i$ denote passenger \emph{request identity} $i$, $H_j$ denote parcel
request identity $j$, and $T_k$ denote taxi identity $k$. These symbols are
resources, not service-vertex labels. Let
\[
 \Q=\{P_1,\ldots,P_{120},H_1,\ldots,H_{120}\}.
\]
For taxi $k$, let $\R_k$ be all feasible complete routes, including the empty
route. For a route $R\in\R_k$, write $\Q(R)$ for the requests it serves and
\[
 b(R)=\sum_{r\in\Q(R)}\operatorname{revenue}_r
       -\sum_{(u,v)\text{ traversed by }R}c(u,v).
\]
The second sum includes the origin-to-first-service leg. A feasible fleet
$\F$ chooses at most one nonempty route per taxi, with pairwise disjoint
served-request sets. Its benefit is $B(\F)=\sum_{R\in\F}b(R)$.

\begin{theorem}\label{thm:optimum}
For this input and these constraints,
\[
 \max_{\F\text{ feasible}}B(\F)=13\,279.
\]
Consequently, a benefit of at least $13\,304$, which would improve the original
best-known value $13\,043$ by at least $2\%$, is impossible.
\end{theorem}

The proof has two distinct tasks. First, obtain a universal upper bound for
\emph{every} feasible fleet. Second, supply a feasible fleet attaining that
bound. Sections~\ref{sec:bound}--\ref{sec:certificate} establish the upper
bound; Section~\ref{sec:lower} supplies the attaining fleet. The enumeration
argument in Section~\ref{sec:enumeration} explains why checking a finite
route bank suffices for the universal claim.

\section{Constructing a valid upper bound}\label{sec:bound}

\subsection{Why independently optimizing every taxi is insufficient}

One immediate bound is
\[
 B(\F)\leq \sum_{k=1}^{20}\max_{R\in\R_k}b(R)=16\,216.
\]
This is valid because every chosen route is bounded by its taxi's best route.
It is loose because different taxis' best routes can serve the same request.
To account for that competition, assign each request a nonnegative weight
$p_r$. For any such weights, define the residual taxi value
\[
 a_k^{(0)}=\max\left\{0,\max_{R\in\R_k}
       \left[b(R)-\sum_{r\in\Q(R)}p_r\right]\right\}.
\]
Every route obeys $b(R)\leq a_k^{(0)}+\sum_{r\in\Q(R)}p_r$. Summing over a
feasible fleet charges each taxi and request at most once, giving
\[
 B(\F)\leq \sum_k a_k^{(0)}+\sum_r p_r.
\]
Optimizing these weights for the complete route bank gives the fractional
route-master bound $13\,311\frac{2}{3}$. This already limits every integer
fleet, but does not yet prove $13\,279$. Four additional conflict charges
close the remaining gap.

\subsection{A three-resource conflict that can be charged only once}

A route $R\in\R_k$ uses the resource set
\[
 S(R)=\{T_k\}\cup\Q(R).
\]
Across a feasible fleet, every resource in
$\{T_1,\ldots,T_{20}\}\cup\Q$ appears in at most one selected route.
For any set $G$ of three distinct resources, define
\[
 I_G(R)=\ind\{|S(R)\cap G|\geq 2\}.
\]
Thus $I_G(R)=1$ if the route uses at least two resources from $G$.

\begin{lemma}\label{lem:triple}
For every feasible fleet and every three-resource set $G$,
\[
 \sum_{R\in\F}I_G(R)\leq 1.
\]
\end{lemma}
\begin{proof}
Suppose two selected routes both triggered $G$. Each would use at least two
of its three resources. Two subsets of a three-element set, each containing
at least two elements, must intersect. The two routes would therefore use
the same taxi identity or the same request identity, contradicting fleet
feasibility. The argument also applies when one route uses all three resources.
\end{proof}

For example, let $G=\{H_{55},H_{65},H_{90}\}$. A route serving parcels $55$
and $65$ and a second route serving parcels $65$ and $90$ cannot coexist:
both serve parcel $65$. Any other pair of routes each serving two of these
three parcels has the same unavoidable overlap. The lemma also permits a
taxi identity in $G$, as used by the certificate below.

\subsection{The covering inequality and its fleet-wide sum}

Choose nonnegative request weights $p_r$ and nonnegative conflict weights
$q_j$ for three-resource groups $G_1,\ldots,G_J$. Define
\begin{equation}\label{eq:residual}
 a_k=\max\left\{0,\max_{R\in\R_k}
 \left[b(R)-\sum_{r\in\Q(R)}p_r-\sum_{j=1}^Jq_j I_{G_j}(R)\right]\right\}.
\end{equation}
By the definition of a maximum, \emph{every} feasible route of taxi $k$ obeys
\begin{equation}\label{eq:cover}
 b(R)\leq a_k+\sum_{r\in\Q(R)}p_r+\sum_{j=1}^Jq_j I_{G_j}(R).
\end{equation}
Let $z_k$ indicate whether the fleet uses taxi $k$, let $y_r$ indicate whether
it serves request $r$, and put $h_j=\sum_{R\in\F}I_{G_j}(R)$. Summing
\eqref{eq:cover} over its selected routes gives
\begin{align}
 B(\F)
 &\leq\sum_{k=1}^{20}a_k z_k+\sum_{r\in\Q}p_r y_r+\sum_{j=1}^Jq_j h_j
 \notag\\
 &\leq\sum_{k=1}^{20}a_k+\sum_{r\in\Q}p_r+\sum_{j=1}^Jq_j=:U.\label{eq:upper}
\end{align}
The last step uses $z_k\leq1$, $y_r\leq1$, Lemma~\ref{lem:triple}, and the
nonnegativity of every weight. Unused resources merely leave some terms
uncharged. Equation~\eqref{eq:upper} is valid for any feasible fleet, including
one that an ALNS run has never visited.

\subsection{How the weights were obtained}

The weights can be found by solving the following linear program over a
complete route bank $\mathcal C$:
\begin{align*}
 \min\quad&\sum_k a_k+\sum_r p_r+\sum_j q_j\\
 \text{subject to}\quad&
 a_{k(c)}+\sum_{r\in\Q(c)}p_r+\sum_jq_j I_{G_j}(c)\geq b(c)
       &&(c\in\mathcal C),\\
 &a_k,p_r,q_j\geq0.
\end{align*}
Every feasible set of weights already provides an upper bound; optimality
of this auxiliary LP is unnecessary for verifying that bound. During
weight discovery, the complete route LP initially had value
$13\,311\frac23$. Adding seven valid three-resource inequalities lowered it
to $13\,279\frac23$; an eighth lowered it to $13\,279$. Only four of their
weights are positive in the final certificate. The weights below are integers,
and all route-cover inequalities were subsequently checked exactly.

Equivalently, the discovery LP is the dual of the fractional route-selection
model with the conflict inequalities added. The corresponding integer master
selects columns $x_c\in\{0,1\}$:
\begin{align*}
 \max\quad&\sum_{c\in\mathcal C}b(c)x_c\\
 \text{subject to}\quad&\sum_{c:k(c)=k}x_c\leq1 &&(k=1,\ldots,20),\\
 &\sum_{c:r\in\Q(c)}x_c\leq1 &&(r\in\Q),\\
 &x_c\in\{0,1\} &&(c\in\mathcal C).
\end{align*}
The group rows $\sum_c I_{G_j}(c)x_c\leq1$ are valid for this integer model
by Lemma~\ref{lem:triple}. Their role is to exclude fractional combinations
that need not satisfy the two-of-three conflict. Independent CPLEX and HiGHS
solves of the complete integer master also returned $13\,279$. The explicit
covering proof above does not require trusting their reported optimality gap.

\section{The complete integer weight certificate}\label{sec:certificate}

All indices in the following tables are one-based. Every request weight not
listed is exactly zero. The tables give all weights needed to evaluate
\eqref{eq:cover}; there are no omitted positive charges.

\subsection{Taxi weights}
\begin{center}
\begin{tabular}{rr@{\qquad}rr}
\toprule
Taxi $k$ & $a_k$ & Taxi $k$ & $a_k$\\
\midrule
1 & 483 & 11 & 476 \\
2 & 546 & 12 & 576 \\
3 & 538 & 13 & 423 \\
4 & 552 & 14 & 552 \\
5 & 526 & 15 & 711 \\
6 & 519 & 16 & 564 \\
7 & 459 & 17 & 445 \\
8 & 672 & 18 & 425 \\
9 & 529 & 19 & 405 \\
10 & 619 & 20 & 543 \\
\bottomrule
\end{tabular}
\end{center}
\[
 \sum_{k=1}^{20}a_k=10\,563.
\]

\subsection{Nonzero request weights}
The six nonzero passenger weights are:
\begin{center}
\begin{tabular}{rr@{\qquad}rr@{\qquad}rr}
\toprule
Request & Weight & Request & Weight & Request & Weight\\
\midrule
$P_{4}$ & 34 & $P_{7}$ & 2 & $P_{20}$ & 201 \\
$P_{32}$ & 164 & $P_{51}$ & 46 & $P_{114}$ & 41 \\
\bottomrule
\end{tabular}
\end{center}
The twenty-five nonzero parcel weights are:
\begin{center}
\begin{tabular}{rr@{\qquad}rr@{\qquad}rr}
\toprule
Request & Weight & Request & Weight & Request & Weight\\
\midrule
$H_{2}$ & 71 & $H_{12}$ & 89 & $H_{14}$ & 7 \\
$H_{17}$ & 1 & $H_{19}$ & 55 & $H_{20}$ & 33 \\
$H_{28}$ & 225 & $H_{32}$ & 37 & $H_{40}$ & 44 \\
$H_{43}$ & 205 & $H_{49}$ & 214 & $H_{52}$ & 105 \\
$H_{53}$ & 171 & $H_{55}$ & 50 & $H_{58}$ & 76 \\
$H_{65}$ & 38 & $H_{70}$ & 85 & $H_{76}$ & 26 \\
$H_{83}$ & 12 & $H_{87}$ & 142 & $H_{89}$ & 33 \\
$H_{90}$ & 222 & $H_{101}$ & 21 & $H_{110}$ & 34 \\
$H_{115}$ & 63 &  &  &  &  \\
\bottomrule
\end{tabular}
\end{center}
The passenger weights sum to $488$ and the parcel weights to $2\,059$;
therefore
\[
 \sum_{r\in\Q}p_r=2\,547.
\]

\subsection{Three-resource conflict weights}
\begin{center}
\begin{tabular}{cl r}
\toprule
Group & Three resources & $q_j$\\
\midrule
$G_1$ & $\{H_{55}, H_{65}, H_{90}\}$ & 103 \\
$G_2$ & $\{P_{61}, H_{53}, H_{120}\}$ & 43 \\
$G_3$ & $\{T_{6}, P_{61}, H_{53}\}$ & 21 \\
$G_4$ & $\{P_{7}, H_{39}, H_{89}\}$ & 2 \\
\bottomrule
\end{tabular}
\end{center}
Hence $\sum_j q_j=103+43+21+2=169$. Substitution into
\eqref{eq:upper} gives the claimed universal upper bound:
\begin{equation}\label{eq:numerical}
 \boxed{B(\F)\leq U=10\,563+2\,547+169=13\,279.}
\end{equation}

\subsection{A concrete route-level calculation}

Consider the selected route for taxi $7$, whose official service-vertex
sequence is
\[
 411,\ 439,\ 412,\ 389,\ 197,\ 198,\ 390,\ 440.
\]
It serves $P_{89}$ and $H_{65},H_{76},H_{90}$. Its revenue is $1\,178$, its
travel cost is $330$, and its benefit is $848$. The resource set contains two
elements of $G_1=\{H_{55},H_{65},H_{90}\}$ and triggers no other weighted
group. Its covering inequality is therefore
\[
 \underbrace{848}_{b(R)}
 \leq \underbrace{459}_{a_7}
    +\underbrace{0+38+26+222}_{p_{P_{89}}+p_{H_{65}}+p_{H_{76}}+p_{H_{90}}}
    +\underbrace{103}_{q_1}
 =848.
\]
This example shows how one route is covered. To obtain a universal bound,
the same calculation must hold for every possible route, not just the
selected routes. The exhaustive bank check verified it for all $21\,461$
retained columns: the minimum covering slack was zero and $60$ columns were
tight. It also checked that every $a_k$ equals the actual maximum residual
in \eqref{eq:residual}. The next section justifies the bank's completeness.

\section{Why the route bank covers every feasible route}\label{sec:enumeration}

\subsection{Exhaustive forward enumeration}

For each taxi, the enumerator starts at its own origin with time, parcel
load, accumulated revenue, and travel cost equal to zero. A recursive state
contains the current vertex, earliest departure time, load, picked-request
set, onboard-parcel set, and service sequence. It explores each of the
following legal extensions: deliver any onboard parcel; pick up any unused
parcel that fits capacity; or pick up any unused passenger and immediately
visit that passenger's delivery as an indivisible two-vertex extension.

For a next endpoint $v$ with window $[e_v,\ell_v]$, previous departure $d_u$,
and service duration $s_v$, the transition computes
\[
 \operatorname{arrival}_v=d_u+t(u,v),\qquad
 d_v=\max\{\operatorname{arrival}_v,e_v\}+s_v.
\]
It rejects the transition if arrival exceeds $\ell_v$ or the new load exceeds
capacity. A passenger delivery has no latest bound and follows immediately
after the pickup, with its own travel and service time. Revenue is credited
once per pickup; a route is recorded only when its onboard-parcel set is
empty. At that point every picked request has been completed.

The enumerator continues from every feasible state, including closed routes
and prefixes with nonpositive benefit. It has no beam-width limit, time
limit, route-length limit, or heuristic profitability pruning. Each request
can be picked only once, so every branch has at most $480$ service vertices
and the search terminates. The longest expanded sequence observed in this
instance had $11$ vertices; this was an outcome, not an imposed cutoff.
The cheap test ``current time exceeds a pickup's latest time'' is safe because
travel times are nonnegative. No triangle-inequality pruning is used.

\begin{lemma}\label{lem:earliest}
For a fixed structurally valid service sequence, if any timing is feasible,
its earliest-start timing is feasible.
\end{lemma}
\begin{proof}
Both timings begin at time zero. Suppose the earliest construction leaves
the current vertex no later than another feasible schedule. Adding the same
travel time gives an arrival no later than that schedule's next arrival.
The map $a\mapsto\max\{a,e_v\}+s_v$ is nondecreasing, so earliest departure
from the next vertex is also no later. By induction, all earliest arrivals
are no later than feasible arrivals and therefore meet every latest bound.
Waiting enforces every earliest-service requirement. Load and adjacency
depend only on the sequence, so they are unchanged. This argument uses
static travel times, not a triangle inequality.
\end{proof}

\begin{lemma}\label{lem:complete}
Every feasible complete route appears in the uncompressed enumeration.
\end{lemma}
\begin{proof}
Take any feasible route and examine it from its first service vertex. The
next legal action is either a parcel pickup, an onboard parcel delivery, or
a passenger pickup followed by its adjacent delivery. The enumerator tries
every unused pickup and every onboard delivery, so it includes this action.
The capacity check accepts it because the route is feasible. The timing
check accepts the earliest schedule by Lemma~\ref{lem:earliest}. Repeating
this argument reproduces the entire service sequence. At its final vertex,
all parcels have been delivered, so the route is considered for recording.
No negative-profit prefix or route-length cutoff can eliminate this path.
\end{proof}

\subsection{Why retaining one route per taxi and request set is safe}

For each taxi and served-request set, the bank retains only the complete
positive-benefit route with greatest benefit. This replacement happens when
recording columns; it does not prune recursive states. Consider a discarded
positive route $R$ and its retained representative $\widehat R$. They use the
same taxi and requests, so replacing $R$ by $\widehat R$ in a fleet preserves
all inter-route constraints and cannot reduce its benefit. Each route has
already been checked for its internal time windows, capacity, and adjacency;
no terminal time is shared between taxis.

The same replacement also preserves $S(R)$ and every conflict indicator.
Thus the right side of \eqref{eq:cover} is identical for $R$ and
$\widehat R$, while $b(R)\leq b(\widehat R)$. Checking the retained
representative proves the inequality for every discarded ordering of the
same requests. Nonpositive complete routes satisfy it immediately because
all covering weights are nonnegative; an empty taxi contributes zero.
Consequently the finite bank check proves the covering inequality for
\emph{all} feasible complete routes.

\subsection{Evidence that enumeration completed}

The uncapped Python enumeration expanded $427\,881$ recursive states,
including the $20$ roots, and retained $21\,461$ columns. A separately
implemented C++ enumeration agreed on every retained taxi/request-set key
and its maximum benefit. Its run reported zero truncated labels and zero
unfinished taxis. The counts and the exact bank hash are archived; per-taxi
counts are in Appendix~\ref{sec:counts}.

All numeric input tokens for this instance are integers with magnitude at
most $583$. The Python enumeration stores them as double-precision values,
but along a branch there are at most $480$ service events; time, load,
revenue, and cost computations stay well within exactly representable
integer magnitudes. Its $10^{-8}$ comparison tolerance cannot change an
integer feasibility comparison. The final weight and route-benefit check
uses integer arithmetic, with integral benefits explicitly asserted.

The logical trust boundary is explicit: the mathematical upper-bound
argument requires a complete bank under the stated semantics. A certificate
checked against an incomplete heuristic route pool would not prove global
optimality. Lemmas~\ref{lem:earliest}--\ref{lem:complete} justify the exhaustive
algorithm; the completed counts, cross-implementation agreement, and hashes
identify the executed computation. The checks constitute a computational
certificate, not a formal verification of every line of the software.

\section{A feasible fleet attaining the bound}\label{sec:lower}

The official-format fleet is reproduced in Appendix~\ref{sec:solution}.
Independent forward simulation checked every time window, parcel load,
passenger adjacency, complete pickup/delivery pair, and uniqueness of served
requests. It serves $17$ passengers and $46$ parcels. The current C++ full
checker also accepts the solution. Table~\ref{tab:fleet} gives each route's
revenue, cost, benefit, and certificate charges. Here
$P(R)=\sum_{r\in\Q(R)}p_r$ and
$C(R)=\sum_jq_j I_{G_j}(R)$.

\begin{longtable}{rrrrrrr}
\caption{The attaining fleet and its exact covering charges.}\label{tab:fleet}\\
\toprule
Taxi & Revenue & Cost & $b(R)$ & $a_k$ & $P(R)$ & $C(R)$\\
\midrule
\endfirsthead
\toprule
Taxi & Revenue & Cost & $b(R)$ & $a_k$ & $P(R)$ & $C(R)$\\
\midrule
\endhead
1 & 749 & 266 & 483 & 483 & 0 & 0 \\
2 & 1140 & 348 & 792 & 546 & 246 & 0 \\
3 & 915 & 365 & 550 & 538 & 12 & 0 \\
4 & 1125 & 359 & 766 & 552 & 214 & 0 \\
5 & 1088 & 423 & 665 & 526 & 139 & 0 \\
6 & 1286 & 498 & 788 & 519 & 205 & 64 \\
7 & 1178 & 330 & 848 & 459 & 286 & 103 \\
8 & 1008 & 336 & 672 & 672 & 0 & 0 \\
9 & 1116 & 445 & 671 & 529 & 142 & 0 \\
10 & 1225 & 401 & 824 & 619 & 205 & 0 \\
11 & 939 & 417 & 522 & 476 & 46 & 0 \\
12 & 946 & 363 & 583 & 576 & 7 & 0 \\
13 & 859 & 352 & 507 & 423 & 84 & 0 \\
14 & 1048 & 333 & 715 & 552 & 163 & 0 \\
15 & 944 & 233 & 711 & 711 & 0 & 0 \\
16 & 1087 & 367 & 720 & 564 & 156 & 0 \\
17 & 1160 & 354 & 806 & 445 & 361 & 0 \\
18 & 1019 & 392 & 627 & 425 & 202 & 0 \\
19 & 748 & 306 & 442 & 405 & 35 & 2 \\
20 & 960 & 373 & 587 & 543 & 44 & 0 \\
\midrule
Total & 20540 & 7261 & 13279 & 10563 & 2547 & 169 \\
\bottomrule
\end{longtable}
For every row, $b(R)=a_k+P(R)+C(R)$. In particular, summing the independently
computed benefits gives $B(\F^\star)=13\,279$, so the optimum is at least
$13\,279$. Equation~\eqref{eq:numerical} proves it is at most $13\,279$.
This proves Theorem~\ref{thm:optimum}.

For an explicit contradiction, suppose a feasible fleet had benefit greater
than $13\,279$. Its routes would all satisfy \eqref{eq:cover}. Summing those
inequalities, using each taxi/request at most once and
Lemma~\ref{lem:triple}, would force its benefit to be at most $13\,279$.
That contradicts the assumption. Since benefits are integers here, even a
fleet with benefit $13\,280$ is excluded, and therefore so is one with benefit
$13\,304$.

Relative to the original best-known value $13\,043$, the greatest possible
improvement is
\[
 \frac{13\,279-13\,043}{13\,043}\times100\%
 =1.809399677988\ldots\%.
\]
The requested $2\%$ improvement requires
$\lceil1.02\times13\,043\rceil=\lceil13\,303.86\rceil=13\,304$,
which exceeds the proved optimum by $25$. Repeated ALNS runs attaining the
same benefit are useful search evidence, but the covering bound and attaining
fleet are what establish optimality.

\appendix
\section{The attaining official-format solution}\label{sec:solution}

The first line is total benefit. Subsequent line $k$ gives its number of
service vertices and their one-based input labels, omitting the taxi's
implicit origin. This is the exact saved solution whose hash is listed below.
\begin{verbatim}
13279
6 373 403 395 374 396 404
6 461 329 315 462 330 316
4 425 93 94 426
6 413 357 81 82 414 358
6 483 484 489 375 376 490
8 499 133 134 365 500 479 480 366
8 411 439 412 389 197 198 390 440
6 309 455 123 124 456 310
8 377 301 302 378 433 233 234 434
8 385 345 471 227 228 386 472 346
6 121 122 467 135 136 468
6 287 431 235 236 432 288
6 281 369 282 27 28 370
6 299 283 247 248 284 300
6 469 333 223 224 470 334
6 263 399 95 96 400 264
6 297 363 59 60 364 298
8 293 427 428 294 323 83 84 324
4 437 33 34 438
6 339 475 407 408 476 340
\end{verbatim}

\section{Per-taxi exhaustive enumeration counts}\label{sec:counts}

A state means one invocation of the recursive visit function. Passenger
pickup/delivery is a single extension, so its intermediate pickup is not
counted as a separate recursive state. Columns are the maximum-benefit
positive route for each distinct served-request set of that taxi.
\begin{center}
\begin{tabular}{rrr@{\qquad}rrr}
\toprule
Taxi & States & Columns & Taxi & States & Columns\\
\midrule
1 & 29492 & 1032 & 11 & 12985 & 1035 \\
2 & 13537 & 633 & 12 & 32252 & 1184 \\
3 & 10836 & 879 & 13 & 19416 & 1120 \\
4 & 38548 & 1519 & 14 & 27714 & 1656 \\
5 & 25167 & 1116 & 15 & 9190 & 474 \\
6 & 29295 & 1196 & 16 & 25900 & 1331 \\
7 & 19364 & 873 & 17 & 14884 & 1145 \\
8 & 15913 & 678 & 18 & 21053 & 1157 \\
9 & 30115 & 1586 & 19 & 15680 & 915 \\
10 & 10784 & 844 & 20 & 25756 & 1088 \\
\midrule
\multicolumn{4}{r}{Totals:} & $427\,881$ & $21\,461$\\
\bottomrule
\end{tabular}
\end{center}

\section{Reproducibility and artifact identity}\label{sec:reproduce}

Run the commands below from the repository root. The retained audit directory
is \path{results/r240-enhancement-20261004/independent-exact-audit/}; the integer
weights are in
\path{results/r240-enhancement-20261004/algebraic-upper-bound/certificate.json}.
The route-bank JSON contains the taxi, request identities, service sequence,
and exact integral benefit of every retained column.

\subsection{Recheck the bound without an optimizer}
\begin{verbatim}
python3 tests/verify_r240_algebraic_bound.py
\end{verbatim}
The expected result has \texttt{verified=true}, \texttt{upper\_bound=13279},
\texttt{columns=21461}, \texttt{tight\_columns=60},
\texttt{arithmetic="exact integers"}, and \texttt{optimizer\_required=false}.
The verifier first checks the input and bank hashes. It then checks
nonnegative integer weights, three distinct resources per group, coverage
of all columns, each taxi's maximum residual, and the final sum.

For a lower-bound feasibility recheck independent of the C++ evaluator:
\begin{verbatim}
python3 tests/check_r240_independent.py \
  data/probation/r240_p120_m120_s5.txt \
  data/probation/best-known/r240_p120_m120_s5.sol
\end{verbatim}
The expected result is feasible, with benefit $13\,279$, $17$ passengers,
and $46$ parcels. The saved audit also contains the complete per-route
simulation report.

\subsection{Rebuild the complete bank}

Choose an output directory that does not already exist. For example:
\begin{verbatim}
PROOF=results/r240-enhancement-20261004
python3 tests/certify_r240_route_master.py \
  --output /tmp/r240-proof-reenumerated \
  --cpp-bank "$PROOF/beam-route-generation/exhaustive.csv"
\end{verbatim}
This enumerates the Python search tree without a search cap, checks every
retained served-set value against the separate C++ bank, and writes
\texttt{columns.json}, \texttt{all-routes.csv}, \texttt{complete-master.lp},
and \texttt{enumeration.json}. The output bank must have the bank hash below.
The private helpers and result artifacts are retained locally under ignored
\texttt{tests/} and \texttt{results/}; this TeX file has no runtime dependency
on them for typesetting. Reproducing the computational check requires those
files as well as the input.

\subsection{SHA-256 identifiers}

The following hashes identify the input, computational bank, certificate,
and attaining fleet. Each digest is printed in two consecutive 32-character
lines; concatenate the lines without a space or newline.
\noindent\textbf{Instance}\par
\begin{verbatim}
data/probation/r240_p120_m120_s5.txt
93baf239e2d2ccb44cfb2b2a14bb15c4
2804acfc6f8f186cc9f05896e62b86e4
\end{verbatim}

\noindent\textbf{Complete retained route bank}\par
\begin{verbatim}
results/r240-enhancement-20261004/independent-exact-audit/columns.json
5a11b688a68fe22d689defd37239a29f
a934d0bd58d628534351c92fdd54ba1b
\end{verbatim}

\noindent\textbf{Integer weight certificate}\par
\begin{verbatim}
results/r240-enhancement-20261004/algebraic-upper-bound/certificate.json
7a88edd4b54242e81409a4fb4ffe1a3c
d4b35f38cbb07921f98e0327b8f2976d
\end{verbatim}

\noindent\textbf{Attaining official-format solution}\par
\begin{verbatim}
data/probation/best-known/r240_p120_m120_s5.sol
d49bf6d9fffbba6113efb671b86afbd2
e6b0d3b2eae1abad632cca73877110ac
\end{verbatim}

\noindent\textbf{Uncapped Python enumeration source}\par
\begin{verbatim}
tests/certify_r240_route_master.py
0937acfb2ed99d2382ed6f369059fde8
3462a6baafb584e233f2f9d642e30c43
\end{verbatim}

\noindent\textbf{Integer certificate verifier source}\par
\begin{verbatim}
tests/verify_r240_algebraic_bound.py
5964f4e6eabb62e805f29268bcedf471
dd73b69121c21557c46318e25d408b6e
\end{verbatim}

\noindent\textbf{Complete integer master}\par
\begin{verbatim}
results/r240-enhancement-20261004/independent-exact-audit/complete-master.lp
7f6c734fe9e5c6f756141541f3c45ae6
ea88c1c8f71a2b5d52350da50b012bd2
\end{verbatim}

\end{document}
